\hypertarget{writing-a-c-extension-from-scratch}{%
\chapter{Writing a C Extension from
Scratch}\label{writing-a-c-extension-from-scratch}}

True ``Godhood'' involves the ability to extend the Python interpreter
with performance-critical code written in C. This chapter provides a
complete walk-through of creating a high-performance math extension.

\hypertarget{the-anatomy-of-a-c-extension}{%
\subsection{63.1 The Anatomy of a C
Extension}\label{the-anatomy-of-a-c-extension}}

A C extension is a shared library (\passthrough{\lstinline!.so!} or
\passthrough{\lstinline!.pyd!}) that exports an initialization function.

\hypertarget{header-and-types}{%
\subsubsection{1. Header and Types}\label{header-and-types}}

Every extension must include \passthrough{\lstinline!Python.h!}. This
header defines all the \passthrough{\lstinline!PyObject!} structures and
C-API functions.

\begin{lstlisting}[language=C]
#include <Python.h>

// A simple C function to add two numbers
static PyObject* godhood_add(PyObject* self, PyObject* args) {
    long a, b;
    // Parse positional arguments from Python to C types
    if (!PyArg_ParseTuple(args, "ll", &a, &b)) {
        return NULL; // Returns TypeError if parsing fails
    }
    return PyLong_FromLong(a + b); // Convert C long back to Python PyObject
}
\end{lstlisting}

\hypertarget{method-table-and-module-definition}{%
\subsubsection{2. Method Table and Module
Definition}\label{method-table-and-module-definition}}

You must tell Python which functions are exported.

\begin{lstlisting}[language=C]
static PyMethodDef GodhoodMethods[] = {
    {"add",  godhood_add, METH_VARARGS, "Add two numbers in C."},
    {NULL, NULL, 0, NULL}        /* Sentinel */
};

static struct PyModuleDef godhoodmodule = {
    PyModuleDef_HEAD_INIT,
    "godhood",   /* name of module */
    NULL,       /* module documentation */
    -1,         /* size of per-interpreter state of the module, or -1 if the module keeps state in global variables. */
    GodhoodMethods
};
\end{lstlisting}

\hypertarget{initialization-function}{%
\subsubsection{3. Initialization
Function}\label{initialization-function}}

The function name must be
\passthrough{\lstinline!PyInit\_<modulename>!}.

\begin{lstlisting}[language=C]
PyMODINIT_FUNC PyInit_godhood(void) {
    return PyModule_Create(&godhoodmodule);
}
\end{lstlisting}

\hypertarget{compiling-with-setuptools}{%
\subsection{\texorpdfstring{63.2 Compiling with
\texttt{setuptools}}{63.2 Compiling with setuptools}}\label{compiling-with-setuptools}}

You use a \passthrough{\lstinline!setup.py!} file to handle the
platform-specific compilation details.

\begin{lstlisting}[language=Python]
from setuptools import setup, Extension

module = Extension('godhood', sources=['godhood.c'])

setup(name='GodhoodExtension',
      version='1.0',
      description='C extension for high-performance math',
      ext_modules=[module])
\end{lstlisting}

\hypertarget{reference-counting-and-memory-safety}{%
\subsection{63.3 Reference Counting and Memory
Safety}\label{reference-counting-and-memory-safety}}

\textbf{Godhood Warning}: In C, you are responsible for reference
counts. * \textbf{\passthrough{\lstinline!Py\_INCREF(obj)!}}: Increment
count (you are keeping a reference). *
\textbf{\passthrough{\lstinline!Py\_DECREF(obj)!}}: Decrement count (you
are finished with it). * \textbf{Leakage}: Failure to
\passthrough{\lstinline!DECREF!} leads to permanent memory leaks. *
\textbf{Segfaults}: \passthrough{\lstinline!DECREF!}ing an object you
don't own leads to use-after-free crashes.

\begin{center}\rule{0.5\linewidth}{0.5pt}\end{center}
